Let $p_n$ denote the $n$th prime and set
\[
 d_n := p_{n+1}-p_n \qquad (n\ge 1).
\]
Define subsets of the positive integers:
\[
A:=\{n\ge 1: d_{n+1}\ge d_n\},\quad
B:=\{n\ge 1: d_{n+1}\le d_n\},\quad
E:=\{n\ge 1: d_{n+1}=d_n\}.
\]
For a set $S\subseteq\mathbb N$ the \emph{(natural) density} is
\[
\delta(S):=\lim_{N\to\infty}\frac{|S\cap\{1,2,\dots,N\}|}{N}
\]
if the limit exists.

The problem statement (as given) asserts:
\begin{enumerate}[label=(\alph*)]
\item $\delta(A)=\tfrac12$ and $\delta(B)=\tfrac12$;
\item $E$ is infinite (infinitely many $n$ with $d_{n+1}=d_n$);
\item (Erd\H{o}s, 1985) for every $k\ge 1$ there exists $n$ such that
\[d_n=d_{n+1}=\cdots=d_{n+k}.
\]
\end{enumerate}

\medskip
\noindent\textbf{Minimal clarifications / equivalences.}
\begin{itemize}
\item The sets $A$ and $B$ always satisfy $A\cup B=\mathbb N$ and $A\cap B=E$.
\item If both $\delta(A)$ and $\delta(B)$ exist and equal $\tfrac12$, then necessarily $\delta(E)=0$ (because $\delta(A)+\delta(B)=1+\delta(E)$).
\item The condition $d_n=d_{n+1}=\cdots=d_{n+k}$ is equivalent to saying that
$p_n,p_{n+1},\dots,p_{n+k+1}$ form an arithmetic progression (and in fact are \emph{consecutive} primes).
A short proof is given below.
\end{itemize}
